\documentclass[11pt]{article}
\usepackage[margin=0.75in]{geometry}
\usepackage[hidelinks]{hyperref}
\usepackage{parskip}
\pagestyle{empty}
\setlength{\parindent}{0pt}
\setlength{\parskip}{10pt}
\newcommand{\contact}[1]{{\small #1}}

\begin{document}
\begin{center}
{\LARGE \textbf{Yousif Shasha}}\\[4pt]
\contact{Mississauga, ON \enspace|\enspace (647) 887-2009 \enspace|\enspace
\href{mailto:yousif.shasha0@gmail.com}{yousif.shasha0@gmail.com}}\\[3pt]
\contact{\href{https://www.linkedin.com/in/yousif-shasha-177412375/}{LinkedIn}
\enspace|\enspace
\href{https://yousifshasha.github.io/Yousif-Shasha.github.io/}{Portfolio}}
\end{center}
\vspace{-4pt}

September 8, 2026

Hiring Team\\
Hatch\\
Mississauga, ON

\textbf{Re: Water Conveyance Engineering Student -- Mississauga (100526)}

Dear Hatch Hiring Team,

I am a Structural Engineering student at Toronto Metropolitan University applying for the Water Conveyance Engineering Student position in Mississauga. My academic experience in hydraulic modelling, AutoCAD drafting, and surveying provides a foundation for supporting engineering calculations, drawings, and field data collection. I am interested in learning from Hatch's Linear Conveyance team and contributing to municipal water infrastructure projects through their design and construction phases.

In my hydraulic engineering coursework, I built EPANET scenarios to assess how pipe roughness, diameter, demand, and redundancy affect network flow and pressure. Across seven laboratory experiments, I investigated pressurized flow, pipe losses, weirs, and sluice-gate behaviour. I also used HEC-RAS and Excel to compare flood profiles up to the 100-year event and assess velocity and erosion risk. These projects gave me practice evaluating system responses and interpreting results, which I would apply to hydraulic calculations and basic modelling under supervision.

My structural mechanics projects involved preparing dimensioned AutoCAD drawings for aluminum, timber, and composite specimens, performing hand and Excel calculations, and comparing predictions with laboratory measurements. In surveying, I collected total-station elevation, angle, and distance measurements, checked traverse closure, and applied coordinate adjustments. These experiences developed my attention to dimensions, measurement accuracy, and consistency between calculations and physical observations, providing a basis for learning Hatch's drawing and quality-review procedures.

On TMU's Concrete Canoe construction team, I fabricate components, test alternative materials and fabrication methods, and document prototype results while collaborating on design improvements. This work has strengthened my interest in how design decisions affect fabrication and constructability. As a Grocery Supervisor at Groupe Adonis, I also supervise two to six employees, coordinate daily tasks, train new hires, and resolve customer and operational issues. I would bring this communication, organization, and adaptability to working with Hatch's engineers, designers, and project teams.

I am drawn to the opportunity to support watermain and sewer projects that serve communities while developing practical consulting experience. I am eager to learn municipal standards, document control processes, and construction-support tasks such as shop drawing reviews and design clarifications. Thank you for considering my application. I would welcome the opportunity to discuss how my academic foundation and willingness to learn could contribute to your team.

Sincerely,\\[4pt]
Yousif Shasha

\end{document}
